Các nghiên cứu liên quan đến đồ án trải trên ba hướng: quan sát lớp học tự động, khám phá liên kết thời gian và LLM hỗ trợ suy luận đồ thị. Hướng thứ nhất tập trung vào việc giảm công sức mã hóa hành vi thủ công. Wang và cộng sự chứng minh tính khả thi của việc dùng học sâu để nhận diện nhiều hành vi học tập trong lớp học thông minh \,\cite{wang2023smartclassroom}. Yang và cộng sự xây dựng khung dữ liệu và benchmark cho phân tích hành vi học sinh trong điều kiện lớp học thực tế, nơi có đông học sinh, che khuất, chồng lấn hành vi và khác biệt lớn về góc nhìn \,\cite{yang2023yolov7bra}. Các nghiên cứu này cho thấy video lớp học có thể được xử lý tự động để tạo bằng chứng hành vi, nhưng cũng cho thấy môi trường lớp học phức tạp hơn nhiều bối cảnh phát hiện đối tượng thông thường.

Nhiều công trình gần đây cải thiện detector cho hành vi lớp học bằng cách bổ sung mô-đun đặc trưng hoặc attention. Chen và cộng sự đề xuất một biến thể YOLOv8 có Res2Net để xử lý mục tiêu dày đặc và che khuất \,\cite{chen2023improvedyolov8}. Zhang và cộng sự tích hợp Multi-Head Self-Attention và Coordinate Attention vào YOLOv8 nhằm nắm bắt phụ thuộc ngữ cảnh giữa học sinh, bàn học và đồ vật lớp học \,\cite{zhang2025yolov8attention}. Han và cộng sự đề xuất biến thể YOLOv8 có attention động cho nhận diện hành vi trong lớp học \,\cite{han2025wadyolov8}. Ưu điểm chung của các hướng này là nâng cao chất lượng phát hiện hành vi; hạn chế là kết quả thường vẫn dừng ở nhãn, hộp phát hiện, số đếm hoặc tần suất.

Một số nghiên cứu nhấn mạnh rằng hành vi trong lớp không nên được xem như các sự kiện độc lập. Zhang và cộng sự cho rằng mô hình hóa quan hệ tương tác giữa học sinh có thể cung cấp diễn giải hành vi giàu hơn so với việc chỉ tập trung vào từng đối tượng riêng lẻ \,\cite{zhang2024superresolution}. Nhận định này phù hợp với mục tiêu của CausClass: sau khi phát hiện hành vi, hệ thống cần tiếp tục phân tích các mẫu thay đổi theo thời gian thay vì coi detector là điểm kết thúc.

Hướng nghiên cứu thứ hai liên quan đến khám phá liên kết trên chuỗi thời gian. Tổng quan của Runge nhấn mạnh rằng phân tích dữ liệu tuần tự cần xét đến thứ tự thời gian, phụ thuộc có điều kiện và nhiễu ẩn \,\cite{runge2023causalreview}. Truyền thống Granger cung cấp một cách hiểu phù hợp cho dữ liệu quan sát: một biến có liên kết dự báo đến biến khác nếu lịch sử của biến đó giúp cải thiện dự báo biến đích \,\cite{granger1969causality}. Các phương pháp hiện đại như CUTS+ và DyCAST mở rộng hướng này bằng mô hình nơ-ron và đồ thị động để xử lý chuỗi cao chiều, phi tuyến hoặc cấu trúc thay đổi theo thời gian \,\cite{cheng2024cutsplus,cheng2025dycast}. Tuy nhiên, khi áp dụng cho lớp học, các phương pháp này cần một bước trung gian để chuyển video thành chuỗi hành vi có ý nghĩa và có thể truy vết.

AERCA là một công trình gần với nhu cầu của đồ án vì kết hợp autoencoder với khám phá phụ thuộc Granger trên chuỗi thời gian đa biến \,\cite{han2025aerca}. Trong thiết kế gốc, AERCA phục vụ phân tích nguyên nhân gốc của bất thường: encoder ước lượng thành phần ngoại sinh, decoder tái tạo quan sát từ lịch sử và biến ngoại sinh dưới cấu trúc phụ thuộc học được. Đối với CausClass, AERCA được sử dụng theo hướng thận trọng hơn: mô hình cung cấp tín hiệu cấu trúc và kiểm chứng liên kết dự báo, không được dùng để kết luận nhân quả chắc chắn trong lớp học.

Hướng nghiên cứu thứ ba dùng LLM để hỗ trợ khám phá đồ thị. Các khảo sát gần đây cho thấy LLM có thể tham gia vào nhiều tác vụ nhân quả, bao gồm trích xuất quan hệ, sinh prior, hỗ trợ định hướng cạnh và diễn giải cấu trúc \,\cite{wan2024llmcausal,ma2025llmcausal}. Darvariu và cộng sự cho thấy prior do LLM sinh có thể giảm không gian tìm kiếm và hỗ trợ định hướng cạnh khi bằng chứng quan sát chưa đủ rõ \,\cite{darvariu2024llmpriors}. Roy và cộng sự đề xuất Causal-LLM như một khung kết hợp prompt và dữ liệu để suy luận cấu trúc toàn cục \,\cite{roy2025causalllm}. Hạn chế chung là LLM có thể đề xuất quan hệ hợp lý về mặt ngôn ngữ nhưng không được dữ liệu ủng hộ, nên cần cơ chế kiểm chứng độc lập.

Từ các nghiên cứu tương tự, khoảng trống mà đồ án hướng tới là sự thiếu liên kết giữa ba thành phần: perception thị giác, khám phá phụ thuộc thời gian và suy luận có hướng dẫn bởi tri thức. Nếu chỉ có detector, hệ thống tạo ra nhãn và tần suất nhưng ít khả năng giải thích quan hệ thời gian. Nếu chỉ có mô hình đồ thị, kết quả có thể tách khỏi nguồn gốc video và khó kiểm tra lại. Nếu chỉ có LLM, kết quả có thể hợp lý về mặt diễn đạt nhưng thiếu bằng chứng định lượng. CausClass được xây dựng để nối ba thành phần này trong một pipeline có kiểm chứng.